\documentclass[10pt,a4paper]{amsart}
\usepackage[margin=19mm]{geometry}
\usepackage{amsmath,amssymb,mathtools}
\usepackage[scr=rsfs,cal=euler]{mathalfa}
\usepackage{xfrac}
\usepackage[unicode,hidelinks,bookmarks=false]{hyperref}
\newcommand{\ie}{i.e.\ }
\newcommand{\cf}{cf.\ }
\newcommand{\resp}{respectively\ }
\newcommand{\Subsection}{\subsection}
\providecommand{\nobreakdash}{\nobreak\hbox{-}}
\setlength{\emergencystretch}{2em}
\multlinegap0pt
\usepackage{bm}
\usepackage{bbm}
\usepackage{dsfont}
\usepackage{xspace}
\usepackage{cancel}
\usepackage{mathtools}
\usepackage{amsthm}
\makeatletter
\@ifundefined{theorem}{\newtheorem{theorem}{Theorem}}{}
\@ifundefined{lemma}{\newtheorem{lemma}[theorem]{Lemma}}{}
\makeatother
\newcommand{\Z}{\mathbb{Z}}
\renewcommand{\mod}[1]{\;(\text{mod }#1)}
\title[Bounds on Arithmetic Rainbow Ramsey Multiplicities]{Bounds on Arithmetic Rainbow Ramsey Multiplicities}
\author{Gabriel Elvin, Alexis Gonzales, Alejandro Rodriguez,, Israel Wilbur}
\date{Source: arXiv:2601.05442v1 (CC BY 4.0)}
\begin{document}
\maketitle

\noindent\emph{Source and licence.} Bounds on Arithmetic Rainbow Ramsey Multiplicities. Version arXiv:2601.05442v1 (CC BY 4.0), distributed under the Creative Commons Attribution 4.0 International licence, as recorded in the arXiv licence metadata for this exact version and shown on the abstract page https://arxiv.org/abs/2601.05442v1. Full rights notice: licenses/arxiv-2601.05442-CC-BY-4.0.txt.\par\smallskip

\noindent\textit{Editorial scope.} Counted unit: the Lemma whose source label is \texttt{lem-solns-Zn} in the article (source0.tex lines 365--367, source label \texttt{lem-solns-Zn}), delivered together with the article's own complete original proof of it (source0.tex lines 369--398, one \texttt{proof} environment of 30 lines). No mathematics is written, repaired or re-derived: statement and proof are the article's own text line for line. Required context retained uncounted: thm-nt-1 (lines 353--355); thm-nt-2 (lines 357--359); thm-nt-3 (lines 361--363). Every definition, display and label of those lines is retained unchanged. Omitted from this package and not redistributed: 1--352, 356--356, 360--360, 364--364, 368--368, 399--639. Nothing inside a retained excerpt is omitted or altered: every retained body is delivered exactly as the article's lines are written, so no omission receipt of any kind accompanies this package. Citations: the retained excerpts carry no citation command at all, so no bibliography entry of the article is transcribed and no reference is redistributed. Reference adaptation: none was needed and none was made; every reference inside the delivered unit resolves inside it, so no unresolved reference or citation is produced. Printed slips kept literal: no printed word, symbol, hypothesis or number of the counted statement or of its proof was altered. Layout and class: the article's own class is \texttt{article} with packages including amsmath, amsthm, amssymb, amstext, dsfont, fancyvrb, float, fontenc, graphicx, subfigure, hyperref, tikz, inputenc, geometry; the wrapper is the lane's pinned \texttt{amsart} editorial wrapper at 10pt on A4, which restates the theorem-like environments the retained text needs and adds the article's own macro definitions that the retained text needs (\texttt{\textbackslash Z}, \texttt{\textbackslash mod}), with extra wrapper packages (bm, bbm, dsfont, xspace, cancel, mathtools). The delivered numbering is the unit's own. Assets: no retained span contains a figure environment, a graphics-inclusion command, a PDF-page inclusion, a table or a TikZ picture; no image file is copied into this package, the package declares no asset dependency and no image file is redistributed with it. Licence: version arXiv:2601.05442v1 of this article is distributed under the Creative Commons Attribution 4.0 International licence, as recorded in the arXiv licence metadata for this exact version and shown on the abstract page https://arxiv.org/abs/2601.05442v1. A full rights notice is delivered at licenses/arxiv-2601.05442-CC-BY-4.0.txt. Characters: every retained excerpt is pure ASCII, so the delivered engine, XeLaTeX with its default Latin Modern text font, reports no missing character.
\begin{theorem}\label{thm-nt-1}
The congruence $ax \equiv b \mod{n}$ is solvable in integers if and only if $\gcd(a,n)|b$.
\end{theorem}

\begin{theorem}\label{thm-nt-2}
    Given integers $a$, $b$, and $c$ with $a$ and $b$ not both $0$, there exist integers $x$ and $y$ that satisfy the equation $ax + by  = c$ if and only if $\gcd(a,b)|c$.
\end{theorem}

\begin{theorem}\label{thm-nt-3}
If $ax \equiv b \mod{n}$ has a solution, then there are exactly $\gcd(a,n)$ solutions in $\{0, 1, 2, \ldots, n - 1\}$, the canonical complete residue system modulo $n$.
\end{theorem}

\begin{lemma}\label{lem-solns-Zn}
There are $n^2$ solutions to $ax + by = cz$ over $\Z_n$, assuming $\gcd(a,b,c) = 1$.
\end{lemma}

\begin{proof}
        We wish to count the number of solutions to
            \begin{equation}\label{alex-1}
                ax + by \equiv cz \mod{n}, \qquad 0 \leq x, y, z \leq n - 1.
            \end{equation}
        For a fixed $x, y$, by Theorem \ref{thm-nt-1}, there exists $z$ which satisfies the congruence (\ref{alex-1}) if and only if
            \begin{equation}\label{alex-2}
                \gcd(c,n) | ax + by.
            \end{equation}
        Putting $d = \gcd(c,n)$, (\ref{alex-2}) gives us
            \begin{equation}\label{alex-3}
                \begin{split}
                    ax + by &= pd,\\
                    pd - by &= ax,
                \end{split}
            \end{equation}
        where $p \in \Z$. Now by Theorem \ref{thm-nt-2}, a $y$ which satisfies (\ref{alex-3}) exists if and only if $\gcd(b,d)$ divides $ax$. Note that
        this is actually if and only if $\gcd(b,d) | x$, since $\gcd(b,d)$ and $a$ are relatively prime,
        as we are assuming the original equation $ax + by = cz$ is fully reduced. 
        Thus, there are $n /\gcd(b,d)$
        many $x$s in $\Z_n$ that will be a multiple of $\gcd(b,d)$ and hence yield a solution. 
        Now suppose $x$ yields a solution, i.e. satisfies $\gcd(b,d)|x$. Taking (\ref{alex-3}) modulo $d$ gives us
            \begin{equation}\label{alex-5}
                -by \equiv ax \mod{d}.
            \end{equation}
        Then by Theorem \ref{thm-nt-3} there are $\gcd(b,d)$ many $y$s in $\Z_d$ that will satisfy (\ref{alex-5}). Note, each solution in $\Z_d$ corresponds to $n/d$ solutions in $\Z_n$. Thus, there will be $\gcd(b,d) \cdot (n/d)$ suitable $y$s in $\Z_n$ for each suitable $x$. Therefore, there will be
            $$\left( \frac{n}{\gcd(b,d)} \right) \cdot \left( \frac{\gcd(b,d) \cdot n}{d} \right) = \frac{n^2}{d}$$
        pairs $(x,y)$ which satisfy (\ref{alex-2}). Again by Theorem \ref{thm-nt-3}, for each such pair there will be exactly $\gcd(c,n) = d$ many $z$s in $\Z_n$ that form a solution to (\ref{alex-1}). Therefore, there are
            $(n^2 / d) \cdot d = n^2$ total solutions.
\end{proof}

\end{document}
